% Escreva aqui a solução da questão 8.


\textcolor{red}{ VERSÃO 1 } %Rafaela

Sejam $a, b, c \in \R_+^\ast$. Aplicando o Teorema 38 em $a, b \in \R_+^\ast $, teremos:
  \begin{align*}
            \sqrt{ab} \leq \frac{a+b}{2} & \implies c\cdot\sqrt{ab} \leq c\cdot\frac{a+b}{2} &\text{(multiplica por $c>0$)}\\
            & \implies c\sqrt{ab} \leq \frac{ca+cb}{2}, \tag{I} 
\end{align*}       
 em $a, c \in \R_+^\ast $:
 \begin{align*}
            \sqrt{ac} \leq \frac{a+c}{2} & \implies b\cdot\sqrt{ac} \leq b\cdot\frac{a+c}{2} &\text{(multiplica por $b>0$)}\\
                & \implies b\sqrt{ac} \leq \frac{ba+bc}{2}  \tag{II} 
        \end{align*}      
 e em $b, c \in \R_+^\ast $:
 \begin{align*}
            \sqrt{bc} \leq \frac{b+c}{2} & \implies a\cdot\sqrt{bc} \leq a\cdot\frac{b+c}{2} &\text{(multiplica por $a>0$)}\\
                & \implies a\sqrt{bc} \leq \frac{ab+ac}{2}.  \tag{III} 
        \end{align*}
        
``Somando'' (I), (II) e (III), temos:
  \begin{align*}
            c\sqrt{ab} + b\sqrt{ac} + a\sqrt{bc} &\leq \frac{ca+cb}{2} + \frac{ba+bc}{2} + \frac{ab+ac}{2} \\ 
           \implies c\sqrt{ab} + b\sqrt{ac} + a\sqrt{bc} &\leq \frac{2ab+2ac+2bc}{2} \\
           \implies c \sqrt{ab} + b\sqrt{ac} + a\sqrt{bc} &\leq  ab + ac + bc.
        \end{align*}
Portanto, $ ab + ac + bc \geq c\sqrt{ab} + b\sqrt{ac} + a\sqrt{bc}$ para quaisquer $a, b, c \in \R_+^\ast$.






\bigskip


\textcolor{red}{ VERSÃO 2} %Rafaela 

Sejam $a, b, c \in \R_+^\ast$. 

Aplicando o Teorema 38 em $ab, ac \in \R_+^\ast$, teremos:
  \begin{align*}
            \sqrt{ab\cdot ac} \leq \frac{ab+ac}{2} & \implies a\sqrt{bc} \leq \frac{ab+ac}{2}. & \tag{I}
\end{align*}       
Agora em $ba, bc \in \R_+^\ast$, teremos:
  \begin{align*}
            \sqrt{ba\cdot bc} \leq \frac{ba+bc}{2} & \implies b\sqrt{ac} \leq \frac{ba+bc}{2}. & \tag{II}
\end{align*}

Por fim, em $ca, cb \in \R_+^\ast$, teremos:
  \begin{align*}
            \sqrt{ca\cdot cb} \leq \frac{ca+cb}{2} & \implies c\sqrt{ab} \leq \frac{ca+cb}{2}. & \tag{III}
\end{align*}
        
``Somando'' (I), (II) e (III), temos:
  \begin{align*}
            a\sqrt{bc} + b\sqrt{ac} + c\sqrt{ab} &\leq \frac{ab+ac}{2} + \frac{ba+bc}{2} + \frac{ca+cb}{2} \\ 
           \implies a\sqrt{bc} + b\sqrt{ac} + c\sqrt{ab} &\leq \frac{2ab+2ac+2bc}{2} \\
           \implies a\sqrt{bc} + b\sqrt{ac} + c\sqrt{ab} &\leq  ab + ac + bc.
        \end{align*}
Portanto, $ ab + ac + bc \geq c\sqrt{ab} + b\sqrt{ac} + a\sqrt{bc}$ para quaisquer $a, b, c \in \R_+^\ast$.











%\bigskip

%\textcolor{red}{ VERSÃO 3} % autor




%\bigskip





\bigskip

\textcolor{red}{CHAVE DE PONTUAÇÃO:
\begin{itemize}
    \item  Obtenção das 3 desigualdades via Teorema 38 \emph{(1,0 ponto)};
    \item Manipulação das inequações  \emph{(1,0 ponto)};
    \item  Redação em si (desenho, notação e organização lógica) \emph{(1,0 ponto)}.
\end{itemize} }